\documentclass[10pt]{extarticle}
\usepackage[a4paper,margin=1.5cm]{geometry}
\usepackage[T1]{fontenc}
\usepackage{amsmath,amssymb}
\usepackage{graphicx}
\usepackage[export]{adjustbox}
\usepackage{multicol}
\usepackage{enumitem}
\usepackage{booktabs,array,tabularx}
\setlength{\columnsep}{0.7cm}
\setlength{\parindent}{0pt}
\setlength{\parskip}{3pt}
\pagestyle{empty}
\makeatletter
% figure registry (figdefs.tex is generated by the extract step)
\newcommand{\qfigdef}[4]{\@namedef{qf@file@#1}{#2}\@namedef{qf@w@#1}{#3}\@namedef{qf@wide@#1}{#4}}
\newcommand{\qf@check}[1]{\@ifundefined{qf@file@#1}{\PackageError{score2}{Unknown figure id #1}{}}{}}
\newcommand{\qf@img}[4][]{\includegraphics[width=#2\dimexpr\@nameuse{qf@w@#4}pt\relax,max width=#3,max height=0.45\textheight,#1]{\@nameuse{qf@file@#4}}}
% Full-width questions: finalize wraps every question that contains a wide figure in \qfullbegin...\qfullend,
% so the WHOLE question (stem, figures, options) is set across the page in reading order, then the two columns
% resume. The question's head - its number, stem and first figure (\qfullsplit marks the end of the head) -
% is kept on one page, so a number is never left at a page foot with its figure overleaf; the rest may break
% across pages normally, so no big gaps are left. Without a \qfullsplit, the whole question is kept together
% if it is under 45% of a page. Never write these yourself.
\newif\ifqf@full
\newif\ifqf@boxed
\newsavebox{\qf@box}
\newcommand{\qf@open}{\global\qf@fulltrue\par\global\qf@boxedtrue
  \setbox\qf@box\vbox\bgroup\hsize\textwidth\linewidth\textwidth}
\newcommand{\qf@emit}[1]{\par\egroup\global\qf@boxedfalse
  \ifdim\dimexpr\ht\qf@box+\dp\qf@box\relax<#1\textheight\noindent\usebox{\qf@box}\par\else\unvbox\qf@box\fi}
\newcommand{\qfullsplit}{\ifqf@boxed\qf@emit{0.85}\fi}
\newcommand{\qf@close}{\ifqf@boxed\qf@emit{0.45}\fi\global\qf@fullfalse}
\newcommand{\qfullbegin}{\end{multicols}\qf@open}
\newcommand{\qfullend}{\qf@close\begin{multicols}{2}}
\newcommand{\qfullstart}{\qf@open}
\newcommand{\qfulldone}{\qf@close}
% \qfig[scale]{id}: block figure, centred, directly where it is written (size set by the tool).
% Fallback if a wide figure is not inside a full-width question: break out of the columns in place - no floats.
\newcommand{\qfig}[2][1]{\qf@check{#2}%
  \ifqf@full\qf@blk{#1}{#2}\else
  \ifnum\@nameuse{qf@wide@#2}=1 \expandafter\@firstoftwo\else\expandafter\@secondoftwo\fi
  {\end{multicols}\begin{center}\qf@img{#1}{\textwidth}{#2}\end{center}\begin{multicols}{2}}%
  {\qf@blk{#1}{#2}}\fi}
\newcommand{\qf@blk}[2]{\par\smallskip{\centering\qf@img{#1}{\linewidth}{#2}\par}\smallskip}
% \qfigi[scale]{id}: inline figure (inside option lists / tables / minipages); never breaks columns.
\newcommand{\qfigi}[2][1]{\qf@check{#2}\qf@img[valign=c]{#1}{\linewidth}{#2}}
\makeatother
% structure macros - lint reads these to check completeness and figure placement
\newcommand{\papersec}[1]{\par\bigskip{\centering\Large\bfseries PAPER-#1\par}\medskip}
\newcommand{\partsec}[2]{\par\medskip\noindent{\large\bfseries PART-#1: #2}\par\smallskip}
\newcommand{\sect}[1]{\par\smallskip\noindent\textbf{#1}\par}
\newcommand{\qn}[1]{\par\smallskip\noindent\textbf{#1.}~}
\newcommand{\paraq}[2]{\par\smallskip\noindent\textbf{Paragraph for Questions #1 and #2}\par}
\newenvironment{qoptions}{\begin{itemize}[leftmargin=1.8em,itemsep=1pt,topsep=2pt,label={}]}{\end{itemize}}
\newcommand{\opt}[2]{\item[(#1)] #2}
% figure registry (generated by score2_tools.py finalize)
\qfigdef{p1_phy_s1_q01_1}{Score Advanced - Test 20 (06 Apr 2024)-img/p1_phy_s1_q01_1.png}{90.6}{0}
\qfigdef{p1_phy_s1_q03_1}{Score Advanced - Test 20 (06 Apr 2024)-img/p1_phy_s1_q03_1.png}{105.6}{0}
\qfigdef{p1_phy_s1_q04_1}{Score Advanced - Test 20 (06 Apr 2024)-img/p1_phy_s1_q04_1.png}{119.9}{0}
\qfigdef{p1_phy_s1_q05_1}{Score Advanced - Test 20 (06 Apr 2024)-img/p1_phy_s1_q05_1.png}{100.6}{0}
\qfigdef{p1_phy_s2_q07_1}{Score Advanced - Test 20 (06 Apr 2024)-img/p1_phy_s2_q07_1.png}{49.7}{0}
\qfigdef{p1_phy_s2_q08_1}{Score Advanced - Test 20 (06 Apr 2024)-img/p1_phy_s2_q08_1.png}{99.7}{0}
\qfigdef{p1_phy_s2_q10_1}{Score Advanced - Test 20 (06 Apr 2024)-img/p1_phy_s2_q10_1.png}{103.8}{0}
\qfigdef{p1_phy_s3_q01_1}{Score Advanced - Test 20 (06 Apr 2024)-img/p1_phy_s3_q01_1.png}{114.9}{0}
\qfigdef{p1_phy_s3_q02_1}{Score Advanced - Test 20 (06 Apr 2024)-img/p1_phy_s3_q02_1.png}{207.5}{0}
\qfigdef{p1_phy_s3_q03_1}{Score Advanced - Test 20 (06 Apr 2024)-img/p1_phy_s3_q03_1.png}{110.0}{0}
\qfigdef{p1_phy_s3_q04_1}{Score Advanced - Test 20 (06 Apr 2024)-img/p1_phy_s3_q04_1.png}{77.0}{0}
\qfigdef{p1_phy_s3_q05_1}{Score Advanced - Test 20 (06 Apr 2024)-img/p1_phy_s3_q05_1.png}{95.5}{0}
\qfigdef{p1_phy_s3_q06_1}{Score Advanced - Test 20 (06 Apr 2024)-img/p1_phy_s3_q06_1.png}{66.3}{0}
\qfigdef{p1_phy_s3_q07_1}{Score Advanced - Test 20 (06 Apr 2024)-img/p1_phy_s3_q07_1.png}{114.2}{0}
\qfigdef{p1_phy_s3_q08_1}{Score Advanced - Test 20 (06 Apr 2024)-img/p1_phy_s3_q08_1.png}{104.8}{0}
\qfigdef{p1_chem_s1_q02_A}{Score Advanced - Test 20 (06 Apr 2024)-img/p1_chem_s1_q02_A.png}{69.3}{0}
\qfigdef{p1_chem_s1_q02_B}{Score Advanced - Test 20 (06 Apr 2024)-img/p1_chem_s1_q02_B.png}{79.1}{0}
\qfigdef{p1_chem_s1_q02_C}{Score Advanced - Test 20 (06 Apr 2024)-img/p1_chem_s1_q02_C.png}{88.4}{0}
\qfigdef{p1_chem_s1_q02_D}{Score Advanced - Test 20 (06 Apr 2024)-img/p1_chem_s1_q02_D.png}{88.8}{0}
\qfigdef{p1_chem_s1_q04_1}{Score Advanced - Test 20 (06 Apr 2024)-img/p1_chem_s1_q04_1.png}{289.7}{1}
\qfigdef{p1_chem_s3_q01_1}{Score Advanced - Test 20 (06 Apr 2024)-img/p1_chem_s3_q01_1.png}{259.1}{1}
\qfigdef{p1_chem_s3_q03_1}{Score Advanced - Test 20 (06 Apr 2024)-img/p1_chem_s3_q03_1.png}{245.2}{0}
\qfigdef{p1_chem_s3_q07_1}{Score Advanced - Test 20 (06 Apr 2024)-img/p1_chem_s3_q07_1.png}{131.1}{0}
\raggedcolumns
\makeatletter\@beginparpenalty=10000\makeatother
\newcolumntype{L}[1]{>{\raggedright\arraybackslash}p{#1}}
\begin{document}
\begin{center}{\Large\bfseries Score Advanced --- Test 20 (06 Apr 2024)}\end{center}
\begin{multicols}{2}
\raggedcolumns
\papersec{1}
\partsec{1}{PHYSICS}
\sect{SECTION-I (i)}
\qn{1} Two collars, A of mass $2m$ and B of mass $3m$ are connected by a massless rod with pinjoints so that they are constrained to move along the fixed perpendicular rods on horizontal plane as shown, in the accompanying figure. Another collar C of mass $m$ hits A with speed '$v$' and sticks to it. There is no friction between collars and rods.\par
\qfig{p1_phy_s1_q01_1}
\begin{qoptions}
\opt{A}{The speed of collar C just after the impact is $\dfrac{v\cos^2\alpha}{3}$}
\opt{B}{The speed of collar A just after the impact is $\dfrac{v\sin^2\alpha}{3}$}
\opt{C}{The speed of collar B just after the impact is $\dfrac{v\sin2\alpha}{6}$}
\opt{D}{The speed of collar B just after the impact is $\dfrac{v\cos2\alpha}{3}$}
\end{qoptions}
\qn{2} A non-relativistic particle of mass $m$ is held in a circular orbit around the origin by an attractive force. The potential energy of the force is given by $U(r)=\dfrac{kr^2}{2}$, where $r$ is the distance of the particle from origin and $k$ is a constant. Assuming the Bohr quantization of the angular momentum of the particle, choose the correct ($v\longrightarrow$ velocity of particle, $KE(r)\longrightarrow$ kinetic energy of particle, $E(r)\longrightarrow$ total energy of the particle)
\begin{qoptions}
\opt{A}{$r^2=\dfrac{nh}{2\pi\sqrt{km}}$}
\opt{B}{$v^2=\dfrac{nh}{2\pi\sqrt{km}}$}
\opt{C}{$KE(r)=U(r)$}
\opt{D}{$E(r)=\dfrac{nh}{2\pi}\sqrt{\dfrac{k}{m}}$}
\end{qoptions}
\qn{3} A rectangular frame of negligible resistance with capacitor of capacitance C is placed nearby an infinitely long wire carrying current $I=I_0\sin\left(\omega t+\dfrac{\pi}{6}\right)$, wire and frame being in same plane. Select correct alternative(s):\par
\qfig{p1_phy_s1_q03_1}
\begin{qoptions}
\opt{A}{Maximum charge on capacitor is $\left(\dfrac{4\mu_0i_0\omega\ \ln2}{\pi}\right)$C}
\opt{B}{Charge on capacitor at $t=0$ is $\left(\dfrac{\sqrt3\mu_0i_0\omega\ \ln2}{\pi}\right)$C}
\opt{C}{Current in frame is minimum at $t=\dfrac{\pi}{3\omega}$}
\opt{D}{Maximum current in frame is $\left(\dfrac{2\omega^2\mu_0i_0\ \ln2}{\pi}\right)$C}
\end{qoptions}
\qn{4} In a resistance wire of unit length and unit cross sectional area, resistivity varies as $P=\dfrac{P_0}{1+x}$. This wire is connected across a battery of emf V.\par
\qfig{p1_phy_s1_q04_1}
A and B are two ends of the wire and C is mid- point. Choose the correct options ?
\begin{qoptions}
\opt{A}{current $I=\dfrac{V}{P_0\ln2}$}
\opt{B}{Electric field at A $=\dfrac{V}{\ln2}$}
\opt{C}{Electric field at B $=\dfrac{3V}{2\ln2}$}
\opt{D}{Electric field at C $=\dfrac{2V}{3\ln2}$}
\end{qoptions}
\qn{5} In Young's double slit experiment at time $t=0$, A is central maxima and B is 4$^{th}$ maxima. Separation between slits and screen is increasing at a constant speed $v$. Then for a later time 't'. Choose the correct option.\par
\qfig{p1_phy_s1_q05_1}
\begin{qoptions}
\opt{A}{At $t=\dfrac{x}{v}$, B will be 2$^{nd}$ order maxima}
\opt{B}{At $t=\dfrac{2x}{v}$, B will be 1$^{st}$ maxima}
\opt{C}{At $t=\dfrac{5x}{3v}$, B will be 2$^{nd}$ minima}
\opt{D}{At $t=\dfrac{7x}{v}$, B will be 1$^{st}$ minima}
\end{qoptions}
\qn{6} A space capsule is in a circular orbit of radius $r_C$ about the earth. A short burst of a rocket in the backward direction increases the kinetic energy by 20\% and causes the capsule to enter an elliptical orbit with the point of firing as the perigee; i.e., $r_p=r_C$.\par
Then choose the CORRECT statement(s) for further motion of capsule:- (Assume mass of capsule to remain nearly constant)
\begin{qoptions}
\opt{A}{Total mechanical energy of capsule decreases.}
\opt{B}{Time period of capsule will decreases.}
\opt{C}{Distance of apogee from centre of earth is $r_A=1.5r$}
\opt{D}{If mechanical energy before the event (short burst) is $E_1$ \& final mechanical energy after the event is $E_f$, then $E_f=0.8E_i$.}
\end{qoptions}
\sect{SECTION-I (ii)}
\qn{7} A uniform rod of length L is moving with speed $v$ on a horizontal smooth surface. One end of the rod is suddenly hinged - this process takes a very small time $L/100v$. Match the ratios in list-I with their values in list-II.\par
\qfig[0.75]{p1_phy_s2_q07_1}
{\footnotesize\setlength{\tabcolsep}{1.5pt}\noindent\begin{tabular}{|c|L{3.9cm}|c|L{1.2cm}|}
\hline
 & \multicolumn{1}{c|}{List-I} & & \multicolumn{1}{c|}{List-II}\\ \hline
(P) & Speed of the centre of mass of the rod just after to just before the event. & (1) & 2/3\\ \hline
(Q) & KE of the rod just after to just before & (2) & 2/9\\ \hline
(R) & Average force exerted by the hinge during hinging to (100x force exerted by hinge just after hinging) & (3) & 3/4\\ \hline
(S) & Time taken to turn the rod by one radian after the event to time taken to cover the length of the rod before the event & (4) & 5/6\\ \hline
\end{tabular}}\par
Code:
\begin{qoptions}
\opt{A}{P $\rightarrow$ 4;Q $\rightarrow$ 2;R $\rightarrow$ 1;S $\rightarrow$ 3}
\opt{B}{P $\rightarrow$ 3;Q $\rightarrow$ 2;R $\rightarrow$ 4;S $\rightarrow$ 1}
\opt{C}{P $\rightarrow$ 3;Q $\rightarrow$ 3;R $\rightarrow$ 2;S $\rightarrow$ 1}
\opt{D}{P $\rightarrow$ 1;Q $\rightarrow$ 4;R $\rightarrow$ 3;S $\rightarrow$ 2}
\end{qoptions}
\qn{8} A liquid of density $\rho$ is completely filled in a cubical container as shown in figure. Neglect the atmospheric pressure, match the quantities in List-I with corresponding expression in List-II and select the CORRECT code\par
\qfig{p1_phy_s2_q08_1}
{\footnotesize\setlength{\tabcolsep}{2pt}\noindent\begin{tabular}{|c|L{3.2cm}|c|L{1.4cm}|}
\hline
 & \multicolumn{1}{c|}{List-I} & & \multicolumn{1}{c|}{List-II}\\ \hline
(P) & Force on rectangular Face 1 is & (1) & $\dfrac13a^3\rho g$\\ \hline
(Q) & Force on rectangular Face 2 is & (2) & $\dfrac18a^3\rho g$\\ \hline
(R) & Force on triangular Face 3 is & (3) & $\dfrac16a^3\rho g$\\ \hline
(S) & Force on triangular Face 4 is & (4) & $\dfrac38a^3\rho g$\\ \hline
\end{tabular}}\par
Code:
\begin{qoptions}
\opt{A}{P $\rightarrow$ 3;Q $\rightarrow$ 1;R $\rightarrow$ 2;S $\rightarrow$ 4}
\opt{B}{P $\rightarrow$ 1;Q $\rightarrow$ 3;R $\rightarrow$ 2;S $\rightarrow$ 4}
\opt{C}{P $\rightarrow$ 2;Q $\rightarrow$ 4;R $\rightarrow$ 1;S $\rightarrow$ 3}
\opt{D}{P $\rightarrow$ 2,4;Q $\rightarrow$ 2;R $\rightarrow$ 1;S $\rightarrow$ 3}
\end{qoptions}
\qn{9} The electric field E is measured at a point P(0,0,d) generated due to various charge distributions and the dependence of E on d is found to be different for different charge distributions. List-I contains different relations between E and d. List-II describes different electric charge distributions, along with their locations. Match the functions in List-I with the related charge distributions in List-II.\par
{\footnotesize\setlength{\tabcolsep}{2pt}\noindent\begin{tabular}{|c|L{1.5cm}|c|L{4.2cm}|}
\hline
 & \multicolumn{1}{c|}{List-I} & & \multicolumn{1}{c|}{List-II}\\ \hline
(P) & E is independent of d & (1) & A point charge Q at the origin\\ \hline
(Q) & E $\propto\dfrac1d$ & (2) & A small dipole with point charges Q at $(0,0,\ell)$ and $-Q$ at $(0,0,-\ell)$. Take $2\ell\ll d$.\\ \hline
(R) & E $\propto\dfrac1{d^2}$ & (3) & An infinite line charge coincident with the x-axis, with uniform linear charge density $\lambda$.\\ \hline
(S) & E $\propto\dfrac1{d^3}$ & (4) & Two infinite wires carrying uniform linear Charge density parallel to the x-axis. The One along $(y=0,z=\ell)$ has a charge density $+\lambda$ and the one along $(y=0,z=-\ell)$ has a charge density $-\lambda$. Take $2\ell\ll d$\\ \hline
 & & (5) & Infinite plane charge coincident with the xy-plane with uniform surface charge density\\ \hline
\end{tabular}}\par
Code:
\begin{qoptions}
\opt{A}{P $\rightarrow$ 5;Q $\rightarrow$ 3,4;R $\rightarrow$ 1;S $\rightarrow$ 2}
\opt{B}{P $\rightarrow$ 5;Q $\rightarrow$ 3;R $\rightarrow$ 1,4;S $\rightarrow$ 2}
\opt{C}{P $\rightarrow$ 5;Q $\rightarrow$ 3;R $\rightarrow$ 1,2;S $\rightarrow$ 4}
\opt{D}{P $\rightarrow$ 4;Q $\rightarrow$ 2,3;R $\rightarrow$ 1;S $\rightarrow$ 5}
\end{qoptions}
\qn{10} A block P of mass $m$ is placed on smooth horizontal surface and connected with a spring. Another block Q of same mass is placed on the block P. The two blocks are pulled by a distance A and released. The blocks oscillates without slipping on each other. Then match the matrix.\par
\qfig{p1_phy_s2_q10_1}
{\footnotesize\setlength{\tabcolsep}{2pt}\noindent\begin{tabular}{|c|L{3.7cm}|c|L{1.9cm}|}
\hline
 & \multicolumn{1}{c|}{List-I} & & \multicolumn{1}{c|}{List-II}\\ \hline
(P) & Friction force on the block Q, when both the blocks are coming towards equilibrium position from left & (1) & Towards right\\ \hline
(Q) & Friction force on the block Q when both block going right from equilibrium position (exclude the equilibrium position) & (2) & Towards left\\ \hline
(R) & Friction force will be maximum on Q & (3) & At equilibrium position\\ \hline
(S) & Friction force on P when both blocks are coming towards the equilibrium position from right & (4) & At extreme position\\ \hline
 & & (5) & None of these\\ \hline
\end{tabular}}\par
Code:
\begin{qoptions}
\opt{A}{P $\rightarrow$ 1;Q $\rightarrow$ 2;R $\rightarrow$ 4;S $\rightarrow$ 1}
\opt{B}{P $\rightarrow$ 2;Q $\rightarrow$ 1;R $\rightarrow$ 4;S $\rightarrow$ 5}
\opt{C}{P $\rightarrow$ 1;Q $\rightarrow$ 2;R $\rightarrow$ 3;S $\rightarrow$ 5}
\opt{D}{P $\rightarrow$ 5;Q $\rightarrow$ 2;R $\rightarrow$ 4;S $\rightarrow$ 3}
\end{qoptions}
\sect{SECTION-II}
\qn{1} Figure shows a cyclic process ABCDBEA performed on an ideal gas. If $P_A=2$ atm.,$P_B=4.5$ atm. and $P_C=6$ atm.; $V_E-V_A=0.05$ litre, find the magnitude of work done by the gas (in J) in the complete process in nearest integer.\par
(1 atm, pressure $=1\times10^5$ pa)\par
\qfig{p1_phy_s3_q01_1}
\qn{2} Gold Nucleus($_{79}$Au$^{198}$) can decay into mercury nucleus ($_{80}$Hg$^{198}$) by two decay schemes shown in figure, (i) It can emit a $\beta$ particle ($\beta_1$) and come to ground state by either emitting one $\gamma$ ray ($\gamma_1$)or emitting two $\gamma$ rays ($\gamma_3$ \& $\gamma_2$). (ii) It can emit one $\beta$ particle ($\beta_2$) and come to ground state by emitting $\gamma_2$ ray. Atomic masses: $^{198}$Au $=197.9682$ amu, $^{198}$Hg $=197.9662$ amu, 1 amu $=930$ MeV/c$^2$. The energy levels of the nucleus are shown in figure. The maximum kinetic energy of emitted $\beta_2$ particles in MeV is equal to\par
\qfig{p1_phy_s3_q02_1}
\qn{3} For the given circuit as shown in figure.\par
\qfig{p1_phy_s3_q03_1}
$R_1=(60\pm3)\,\Omega$\quad $R_2=(30\pm6)\,\Omega$\quad $C=(1\pm0.01)\,\mu$F\quad $E=(40\pm4)$ V\par
At $t=0$ capacitor is uncharged. Maximum fractional error in maximum charge stored in capacitor is \underline{\hspace{1.5cm}} .
\qn{4} Seven particles, each of mass $m$ are placed at the seven corners of a cube of side 'a', but one corner is vacant, as shown in figure. If the x-coordinate of the centre of mass of the system is $\dfrac{na}{m}$. Where $n$ and $m$ are least integer values. Then find value of $(m-n)$.\par
\qfig{p1_phy_s3_q04_1}
\qn{5} For a prism angle $\theta=60^\circ$, the refractive indices of the left half and the right half are, respectively, $n_1$ and $n_2(n_2\ge n_1)$ as shown in the figure. The angle of incidence '$i$' is chosen such that the incident light rays will have minimum deviation if $n_1=n_2=n=1.5$. For the case of unequal refractive indices, $n_1=n$ and $n_2=n+\Delta n$ (where $\Delta n\ll n$), the angle of emergence $e=i+\Delta e$. If $\Delta n=2.8\times10^{-3}$ rad then the value of $\Delta e$ in milliradians is equal to (Round off to nearest decimal up to 2 digits)\par
\qfig{p1_phy_s3_q05_1}
\qn{6} The figure shows an oscillatory system (A uniform rod with an elastic string). The system oscillates in vertical plane shown here. The rubber cord is in natural length when rod is vertical. What is the time period of small oscillation (in sec) ($k$ = force constant of rubber cord). Take $k\ell=mg$, $g=\pi^2$, $\ell=6$m.\par
\qfig{p1_phy_s3_q06_1}
\qn{7} A point source and a receiver are stationary. Source sends a wave of frequency 256 Hz, which reaches to receiver directly and also after being reflected from a wall. The wall is moving at speed of 1 m/s. Find beat frequency (in Hz) heard by receiver. Velocity of sound $=321$ m/s\par
\qfig{p1_phy_s3_q07_1}
\qn{8} A particle is moved along a path AB-BC-CD-DE-EF-FA, as shown in figure, in presence of a force $\vec{F}=\left(\alpha y\hat{i}+2\alpha x\hat{j}\right)$ N, where $x$ and $y$ are in meter and $\alpha=-1$ N/m$^{-1}$. The work done on the particle by this force $\vec{F}$ will be\underline{\hspace{1.5cm}} Joule.\par
\qfig{p1_phy_s3_q08_1}
\partsec{2}{CHEMISTRY}
\sect{SECTION-I (i)}
\qn{1} Which of the following is/are correct?
\begin{qoptions}
\opt{A}{Leaching is the method to remove quantitatively Fe$_2$O$_3$ from red bauxite.}
\opt{B}{In the metallurgy of Aluminium, purified Al$_2$O$_3$ is mixed with Na$_3$AlF$_6$ or CaF$_2$ which lowers the melting point of the mixture and also increases conductivity.}
\opt{C}{Chalcocite and cinnabar are sulphide minerals}
\opt{D}{Self reduction is not possible during metallurgy of Mg.}
\end{qoptions}
\qn{2} Which of the following major organic products produced in the given reactions fails to respond positively in Kjeldahl method of Nitrogen estimation?
\begin{qoptions}
\opt{A}{\qfigi{p1_chem_s1_q02_A}}
\opt{B}{\qfigi{p1_chem_s1_q02_B}}
\opt{C}{\qfigi{p1_chem_s1_q02_C}}
\opt{D}{\qfigi{p1_chem_s1_q02_D}}
\end{qoptions}
\qn{3} The CORRECT statement(s) about the Group-I metals is/are
\begin{qoptions}
\opt{A}{Reactivity of Group-I metals towards water decreases down the group}
\opt{B}{Among the Group-I metals, Li spontaneously reacts with N$_2$ to give a coloured crystalline compound.}
\opt{C}{All the Group-I metal halides are high-melting colourless crystalline solids}
\opt{D}{KCl is more soluble in H$_2$O than NaCl}
\end{qoptions}
\qfullbegin
\qn{4} Correct statement(s) regarding the following reaction sequence is(are) [Consider only major organic product in every step]\par
\qfig{p1_chem_s1_q04_1}
\qfullsplit
\begin{qoptions}
\opt{A}{A on treatment with CHCl$_3$/KOH gives a foul-smelling compound.}
\opt{B}{A is basic while B is amphoteric.}
\opt{C}{C is more reactive than B towards alkaline hydrolysis.}
\opt{D}{F is less reactive than E towards nucleophilic aromatic substitution.}
\end{qoptions}
\qfullend
\qn{5} 10 mole of liquid 'A' and 10 mole of liquid 'B' is mixed in a cylindrical vessel containing a piston arrangement. Initially a pressure of 2 atm is maintained on the solution. Now, the piston is raised slowly and isothermally. Assuming A and B to be completely miscible and forming an ideal solution.\par
$P_A^0=0.6$ atm $P_B^0=0.9$ atm, then choose incorrect answers (Given : $\sqrt{2425}=49.25$, $\sqrt3=1.73$, $\sqrt2=1.41$)
\begin{qoptions}
\opt{A}{The pressure below which the evaporation of liquid solution will start is: 1.75 atm}
\opt{B}{If $\dfrac14$th of the total moles of liquid solution taken initially is converted to vapour, then moles of B in the vapour will be : 3.825}
\opt{C}{The pressure below which the evaporation of liquid solution will start is : 0.75 atm}
\opt{D}{If $\dfrac12$ of the total moles of liquid solution taken initially is converted to vapour, then moles of A in the vapour will be: 2.5}
\end{qoptions}
\qn{6} For the following reaction $2X+Y\xrightarrow{k}P$. The rate of reaction is $\dfrac{d[P]}{dt}=k[X]$.\par
Two moles of X are mixed with one mole of Y to make 1.0 L of solution. At 50 s, 0.5 mole of Y is left in the reaction mixture. The correct statement(s) about the reaction is/are (Use: $\ln2=0.693$)
\begin{qoptions}
\opt{A}{The rate constant, k, of the reaction is $13.86\times10^{-4}$ s$^{-1}$}
\opt{B}{Half-life of X is 50 s}
\opt{C}{At 50s, $-\dfrac{d[X]}{dt}=13.86\times10^{-3}$ mol L$^{-1}$ s$^{-1}$}
\opt{D}{At, 100s, $-\dfrac{d[Y]}{dt}=3.465\times10^{-3}$ mol L$^{-1}$ s$^{-1}$}
\end{qoptions}
\sect{SECTION-I (ii)}
\qn{7} Match the compounds in LIST-I with the observations in LIST-II, and choose the correct option:\par
{\footnotesize\setlength{\tabcolsep}{2pt}\noindent\begin{tabular}{|c|L{1.7cm}|c|L{4.6cm}|}
\hline
 & \multicolumn{1}{c|}{List-I} & & \multicolumn{1}{c|}{List-II}\\ \hline
(P) & Aniline & (1) & Sodium fusion extract of the compound on boiling with FeSO$_4$, followed by acidification with conc. H$_2$SO$_4$, gives Prussian blue colour.\\ \hline
(Q) & o-Cresol & (2) & Sodium fusion extract of the compound on treatment with sodium nitroprusside gives blood red colour.\\ \hline
(R) & Cysteine & (3) & Addition of the compound to a saturated solution of NaHCO$_3$ results in effervescence.\\ \hline
(S) & Caprolactam & (4) & The compound reacts with bromine water to give a white precipitate.\\ \hline
 & & (5) & Treating the compound with neutral FeCl$_3$ solution produces violet colour.\\ \hline
\end{tabular}}\par
Code:
\begin{qoptions}
\opt{A}{P $\rightarrow$ 1,2;Q $\rightarrow$ 4;R $\rightarrow$ 2,3;S $\rightarrow$ 1}
\opt{B}{P $\rightarrow$ 1;Q $\rightarrow$ 3,4;R $\rightarrow$ 3;S $\rightarrow$ 2,4}
\opt{C}{P $\rightarrow$ 2,4;Q $\rightarrow$ 1,5;R $\rightarrow$ 1;S $\rightarrow$ 4}
\opt{D}{P $\rightarrow$ 1,4;Q $\rightarrow$ 5;R $\rightarrow$ 2,3;S $\rightarrow$ 1}
\end{qoptions}
\qn{8} Match the coloured inorganic compounds in List-I with the appropriate phenomenon responsible for the colour in List-II:\par
{\footnotesize\setlength{\tabcolsep}{2pt}\noindent\begin{tabular}{|c|L{3.3cm}|c|L{3.0cm}|}
\hline
 & \multicolumn{1}{c|}{List-I} & & \multicolumn{1}{c|}{List-II}\\ \hline
(P) & PbI$_2$ (Yellow) & (1) & Polarisation\\ \hline
(Q) & (NH$_4$)$_2$[Ce(NO$_3$)$_6$] (Orange-red) & (2) & Ligand $\rightarrow$ Metal charge transfer\\ \hline
(R) & CuSO$_4$.5H$_2$O (Blue) & (3) & f $\rightarrow$ f transition\\ \hline
(S) & KMnO$_4$ (Purple) & (4) & Metal $\rightarrow$ Ligand charge transfer\\ \hline
 & & (5) & d $\rightarrow$ d transition\\ \hline
\end{tabular}}\par
Codes:
\begin{qoptions}
\opt{A}{P $\rightarrow$ 1;Q $\rightarrow$ 2;R $\rightarrow$ 5;S $\rightarrow$ 2}
\opt{B}{P $\rightarrow$ 4;Q $\rightarrow$ 5;R $\rightarrow$ 4;S $\rightarrow$ 2}
\opt{C}{P $\rightarrow$ 4;Q $\rightarrow$ 2;R $\rightarrow$ 5;S $\rightarrow$ 4}
\opt{D}{P $\rightarrow$ 1;Q $\rightarrow$ 4;R $\rightarrow$ 5;S $\rightarrow$ 5}
\end{qoptions}
\qn{9} For the reaction: A $\rightarrow$ product, $t_{x/y}$ represents the times in which $x/y$ fraction of reactant is converted into products.\par
{\footnotesize\setlength{\tabcolsep}{2pt}\noindent\begin{tabular}{|c|L{1.0cm}|c|L{5.2cm}|}
\hline
 & \multicolumn{1}{c|}{List-I} & & \multicolumn{1}{c|}{List-II}\\ \hline
(P) & $t_{5/9}$ & (1) & Equal to 54s, if $t_{1/3}$ is 18s in case of first-order reaction\\ \hline
(Q) & $t_{19/27}$ & (2) & Equal to 32s, if $t_{1/4}$ is 16s in case of first-order reaction\\ \hline
(R) & $t_{7/8}$ & (3) & Equal to 56s, if $t_{1/3}$ is 4s in case of second-order reaction\\ \hline
(S) & $t_{7/16}$ & (4) & Equal to 30s, if $t_{1/3}$ is 18s in case of zero-order reaction.\\ \hline
 & & (5) & Equal to 28s, if $t_{1/2}$ is 16s in case of zero-order reaction.\\ \hline
\end{tabular}}\par
Codes:
\begin{qoptions}
\opt{A}{P $\rightarrow$ 4;Q $\rightarrow$ 1;R $\rightarrow$ 5;S $\rightarrow$ 2}
\opt{B}{P $\rightarrow$ 4;Q $\rightarrow$ 1;R $\rightarrow$ 2;S $\rightarrow$ 5}
\opt{C}{P $\rightarrow$ 4;Q $\rightarrow$ 2;R $\rightarrow$ 5;S $\rightarrow$ 3}
\opt{D}{P $\rightarrow$ 1;Q $\rightarrow$ 4;R $\rightarrow$ 2;S $\rightarrow$ 5}
\end{qoptions}
\qn{10} Match the columns (assume the unit cell edge length equal to $a$).\par
{\footnotesize\setlength{\tabcolsep}{2pt}\noindent\begin{tabular}{|c|L{1.0cm}|c|L{5.2cm}|}
\hline
 & \multicolumn{1}{c|}{List-I} & & \multicolumn{1}{c|}{List-II}\\ \hline
(P) & $a$ & (1) & Two times the smallest distance between two similarly charged ions in fluorite structure.\\ \hline
(Q) & $\dfrac{a}{\sqrt2}$ & (2) & Distance between the next nearest ions in NaCl structure.\\ \hline
(R) & $\sqrt3a$ & (3) & Distance between the next nearest ions in CsCl structure.\\ \hline
(S) & $\sqrt2a$ & (4) & Four times the smallest distance of tetrahedral voids from the corner of FCC unit cell\\ \hline
 & & (5) & Smallest distance between the nearest similarly charged ions in antifluorite structure.\\ \hline
\end{tabular}}\par
Codes:
\begin{qoptions}
\opt{A}{P $\rightarrow$ 1,5;Q $\rightarrow$ 3;R $\rightarrow$ 2;S $\rightarrow$ 4}
\opt{B}{P $\rightarrow$ 1,5;Q $\rightarrow$ 3;R $\rightarrow$ 4;S $\rightarrow$ 2}
\opt{C}{P $\rightarrow$ 2,5;Q $\rightarrow$ 2;R $\rightarrow$ 3;S $\rightarrow$ 4}
\opt{D}{P $\rightarrow$ 2,5;Q $\rightarrow$ 1;R $\rightarrow$ 2;S $\rightarrow$ 4}
\end{qoptions}
\qfullbegin
\sect{SECTION-II}
\qn{1} Lithium borohydride (LiBH$_4$) is a useful reducing agent, as it is more selective than LiAlH$_4$ but less selective than NaBH$_4$. Unlike the later, LiBH$_4$ will reduce esters. However, it will not reduce amides. This selectivity is frequency employed by synthetic organic chemists. In the following reaction sequence, 50.6gm of product (3) is obtained with yield of 65\% and 75\% respectively in 1$^{st}$ and 2$^{nd}$ steps. Find the number of moles of reactant used in this reaction? (C = 12, O = 16, N = 14, H = 1)\par
\qfig{p1_chem_s3_q01_1}
\qfullend
\qn{2} Consider the following species:\par
(i) IF$_5$\quad (ii) ICl$_4^-$\quad (iii) XeO$_2$F$_2$\par
(iv) NH$_2^-$\quad (v) BCl$_3$\quad (vi) BeCl$_2$ (monomer)\par
(vii) AsCl$_4^+$\quad (viii) B(OH)$_3$\quad (ix) NO$_2^-$\quad (x) ClO$_2^+$\par
Calculate value of $(X+Y-Z)$, where\par
$X$ = total number of compounds in which the central atom used their all three p-orbitals in hybridisation respectively.\par
$Y$ = total number of compounds in which the central atom used only two p-orbitals in hybridisation respectively.\par
$Z$ = total number of compounds in which the central atom used only one p-orbital in hybridisation respectively.
\qn{3} Analyse the following reaction sequence.\par
\qfig{p1_chem_s3_q03_1}
If $X$ is the number of alpha hydrogens in the product [Z] and $Y$ is the total no of methyl groups present in the product [Z], then find the value of $X+Y$.
\qn{4} If $X$ is the number of compounds which on strong heating gives oxygen as a product then find the value of $X/2$.\par
HgO, MnO$_2$, Pb$_3$O$_4$, KMnO$_4$, BaO$_2$, LiNO$_3$, Ag$_2$CO$_3$
\qn{5} A pentapeptide has $-$COOH group on tyrosine. This produces (Gly), valine (Val), Phenylalanine (Phe), Alanine (Ala) and Tyrosine (Tyr) on complete hydrolysis. For this pentapeptide, the number of possible sequences (primary structures) with $-$NH$_2$ group attached to a chiral centre is \underline{\hspace{1.5cm}}
\qn{6} At 143 K. the reaction of XeF$_4$ with O$_2$F$_2$ produces a xenon compound Y. The total number of lone pair(s) of electrons present in the molecule of Y is\underline{\hspace{1.5cm}}
\qn{7} 10 mols of an ideal gas is participating in the Carnot cycle is shown in the diagram given below.\par
Step-I: Isothermal reversible expansion with absorption of $q_1$ of heat.\par
Step-II: Adiabatic reversible expansion.\par
Step-III: Isothermal reversible compression with rejection of $q_2$ of heat.\par
Step-IV: Adiabatic reversible expansion.\par
\qfig{p1_chem_s3_q07_1}
Calculate the value of $\dfrac{V_2V_4}{V_3V_1}$ in S.I. unit.
\qn{8} Hydrogen peroxide solution (20 ml) reacts with a solution of KMnO$_4$(20 ml) acidified with dil. H$_2$SO$_4$ The same volume of KMnO$_4$ solution is just decolorised by 10 ml of MnSO$_4$ in neutral medium simultaneously forming a dark brown ppt. of hydrated MnO$_2$. The Brown ppt. dissolved in 10 ml of 0.2 M sodium oxalate in the presence of dil. H$_2$SO$_4$. If the Molarity of H$_2$O$_2$ is $P$ then $10P$ is
\partsec{3}{MATHEMATICS}
\sect{SECTION-I (i)}
\qn{1} If $a$, $b$, $x$, $y$ are positive natural numbers and $\dfrac1x+\dfrac1y=1$ then $\dfrac{a^x}{y}+\dfrac{b^y}{x}$ is same as :
\begin{qoptions}
\opt{A}{$\dfrac{xa^y+yb^x}{x+y}$}
\opt{B}{$\dfrac{y^2.a^x+x^2.b^y}{x^2+y^2}$}
\opt{C}{$\dfrac{a^x}{x}+\dfrac{b^y}{y}$}
\opt{D}{$\dfrac{a^y}{x}+\dfrac{b^x}{y}$}
\end{qoptions}
\qn{2} Let $\left(1-x^3\right)^n=\displaystyle\sum_{r=0}^na_rx^r(1-x)^{3n-2r}$ $\forall\,x\in R-\{1\}$ then which of the following statement is/are correct?
\begin{qoptions}
\opt{A}{If $n=7$ then maximum value of $a_r$ is at $r=5$ or 6}
\opt{B}{$\displaystyle\sum_{r=0}^na_r$ is divisible by $2^n$}
\opt{C}{$\dfrac{a_3}{a_2}$ is integer $\forall\ n\in N$}
\opt{D}{$\displaystyle\sum_{r=0}^na_r$ is not divisible by $2^n$}
\end{qoptions}
\qn{3} Let $S(n, k)$ denote the number of ways in which $n$ distinct balls can be put into $k$ identical boxes such that no box remains empty ($k\le n$). Which of the following is/are correct?
\begin{qoptions}
\opt{A}{$S(5, 3)=25$}
\opt{B}{$S(6, 3)=S(5, 3)+S(5, 2)$}
\opt{C}{$S(6, 3)=3S(5, 3)+S(5, 2)$}
\opt{D}{The number of onto function from 6 element set of 3 element set is $3!\ S(6, 3)$}
\end{qoptions}
\qn{4} If $(\alpha, \beta, \gamma)$ is a point on the line $\dfrac{x-1}{2}=\dfrac{y+1}{-3}=z$ at distance $2\sqrt7$ from the point $(1,-1,0)$, then the coefficient of $x^{\alpha+\beta+\gamma}$ in the expansion of $\left(2x^2-\dfrac1x\right)^9$ is divisible by
\begin{qoptions}
\opt{A}{21}
\opt{B}{36}
\opt{C}{28}
\opt{D}{42}
\end{qoptions}
\qn{5} Let position vectors of points A, B, C are $\vec{a},\vec{b},\vec{c}$ respectively where $\vec{c}=x\vec{a}+y\vec{b}$
\begin{qoptions}
\opt{A}{If $x+y=1$, $x>0$, $y<0$, then C lies on line AB but not on line segment AB`}
\opt{B}{If $x>0$, $y>0$, $x+y<1$, then C lies inside the triangle $\Delta$ABO (O is origin)}
\opt{C}{If $x<0$, $y>0$ and $x+y>1$, then on plane containing points O, A, B, the points C and origin (O) lies on opposite side of line AB.}
\opt{D}{If $x\le0$, $y\le0$ and $x+y=-1$, then locus of point C is reflection of line segment AB about origin.}
\end{qoptions}
\qn{6} Let A and B are square matrix of odd order such that $A^T=A^2-B$ and $B^T=B^2-A$. Then which of the following statement is/are correct?
\begin{qoptions}
\opt{A}{$\left|B^2-A^2\right|=0$}
\opt{B}{$|A|\neq|B|$}
\opt{C}{$|A|=|B|$}
\opt{D}{$\left|B^2-A^2\right|\neq0$}
\end{qoptions}
\sect{SECTION-I (ii)}
\qn{7} Match the following:\par
{\footnotesize\setlength{\tabcolsep}{2pt}\noindent\begin{tabular}{|c|L{4.0cm}|c|L{1.2cm}|}
\hline
 & \multicolumn{1}{c|}{List-I} & & \multicolumn{1}{c|}{List-II}\\ \hline
(P) & The equation of tangent of the ellipse $\dfrac{x^2}{2}+y^2=1$ which cuts off equal lengths of intercepts on coordinate axis is $y=\pm x+a$, then $a$ can be equal to & (1) & $-\sqrt3$\\ \hline
(Q) & The normal $y=mx-2am-am^3$ to the parabola $y^2=4ax$ subtends a right angle at the vertex then $m$ can be equal to & (2) & $-\sqrt2$\\ \hline
(R) & The equation of the common tangent to parabola $y^2=4x$ and $x^2=4y$ is $x+y+\dfrac{k}{\sqrt3}=0$, then $k$ is equal to & (3) & $\sqrt2$\\ \hline
(S) & An equation of common tangent to parabola $y^2=8x$ and the hyperbola $3x^2-y^2=3$ is $2x+\dfrac{ky}{\sqrt2}+1=0$, then $k$ can be equal to & (4) & $\sqrt3$\\ \hline
\end{tabular}}\par
Code:
\begin{qoptions}
\opt{A}{P $\rightarrow$ 1,4;Q $\rightarrow$ 2,3;R $\rightarrow$ 4;S $\rightarrow$ 2,3}
\opt{B}{P $\rightarrow$ 1,3;Q $\rightarrow$ 2;R $\rightarrow$ 3;S $\rightarrow$ 1,4}
\opt{C}{P $\rightarrow$ 4;Q $\rightarrow$ 1,3;R $\rightarrow$ 2,3;S $\rightarrow$ 1}
\opt{D}{P $\rightarrow$ 1,4;Q $\rightarrow$ 2,3;R $\rightarrow$ 1;S $\rightarrow$ 3,4}
\end{qoptions}
\qn{8} Match the following:\par
{\footnotesize\setlength{\tabcolsep}{2pt}\noindent\begin{tabular}{|c|L{4.6cm}|c|L{1.4cm}|}
\hline
 & \multicolumn{1}{c|}{List-I} & & \multicolumn{1}{c|}{List-II}\\ \hline
(P) & The number of rational points on the ellipse $\dfrac{x^2}{9}+\dfrac{y^2}{4}=1$ is & (1) & Infinite\\ \hline
(Q) & The number of integral points on the ellipse $\dfrac{x^2}{9}+\dfrac{y^2}{4}=1$ is & (2) & 4\\ \hline
(R) & The number of rational points on the ellipse $\dfrac{x^2}{3}+y^2=1$ is & (3) & 2\\ \hline
(S) & The number of integral points on the ellipse $\dfrac{x^2}{3}+y^2=1$ is & (4) & 0\\ \hline
\end{tabular}}\par
Code:
\begin{qoptions}
\opt{A}{P $\rightarrow$ 1;Q $\rightarrow$ 3;R $\rightarrow$ 2;S $\rightarrow$ 3}
\opt{B}{P $\rightarrow$ 1;Q $\rightarrow$ 2;R $\rightarrow$ 1;S $\rightarrow$ 3}
\opt{C}{P $\rightarrow$ 3;Q $\rightarrow$ 1;R $\rightarrow$ 1;S $\rightarrow$ 3}
\opt{D}{P $\rightarrow$ 1;Q $\rightarrow$ 2;R $\rightarrow$ 1;S $\rightarrow$ 4}
\end{qoptions}
\qn{9} Match the following:\par
{\footnotesize\setlength{\tabcolsep}{2pt}\noindent\begin{tabular}{|c|L{4.6cm}|c|L{1.0cm}|}
\hline
 & \multicolumn{1}{c|}{List-I} & & \multicolumn{1}{c|}{List-II}\\ \hline
(P) & Area bounded by $y=x^2+2$ and $y=2\,|x|-\cos\pi x$ is equal to $A$, then $3A$ equals & (1) & 4\\ \hline
(Q) & The value of \resizebox{\linewidth}{!}{$2\cot\left(\cot^{-1}3+\cot^{-1}7+\cot^{-1}13+\cot^{-1}21\right)$} & (2) & $\dfrac83$\\ \hline
(R) & Tangent and normal at the ends A \& C of focal chord AC of parabola $y^2=4x$ intersect at B and D, then minimum area of ABCD is & (3) & 3\\ \hline
(S) & Number of integral values of ordered pair $(\alpha,\beta)$ for which the area common to $x^2+y^2-2\alpha^2x-2\beta^2y+c=0$ and its image in $x+y=1$ is maximum is & (4) & 8\\ \hline
\end{tabular}}\par
Code:
\begin{qoptions}
\opt{A}{P $\rightarrow$ 4;Q $\rightarrow$ 3;R $\rightarrow$ 2;S $\rightarrow$ 3}
\opt{B}{P $\rightarrow$ 1;Q $\rightarrow$ 2;R $\rightarrow$ 1;S $\rightarrow$ 3}
\opt{C}{P $\rightarrow$ 4;Q $\rightarrow$ 3;R $\rightarrow$ 4;S $\rightarrow$ 1}
\opt{D}{P $\rightarrow$ 1;Q $\rightarrow$ 3;R $\rightarrow$ 4;S $\rightarrow$ 2}
\end{qoptions}
\qn{10} Let $f(x)$ be a real valued function defined by $f(x)=x^2-2\,|x|$ and\par
$g(x)=\begin{cases}\text{minimum }\{f(t):-2\le t\le x\}&x\in[-2,0)\\[2pt]\text{maximum }\{f(t):0\le t\le x\}&x\in[0,3]\end{cases}$\par
{\footnotesize\setlength{\tabcolsep}{2pt}\noindent\begin{tabular}{|c|L{4.6cm}|c|L{1.0cm}|}
\hline
 & \multicolumn{1}{c|}{List-I} & & \multicolumn{1}{c|}{List-II}\\ \hline
(P) & $g(x)$ is not continuous at $x$ equal to & (1) & $-2$\\ \hline
(Q) & $g(x)$ is not derivable at $x$ equal to & (2) & 0\\ \hline
(R) & Number of points of local maxima of $g(x)$ is equal to & (3) & 1\\ \hline
(S) & Absolute maximum value of $g(x)$ is equal to & (4) & 2\\ \hline
 & & (5) & 3\\ \hline
\end{tabular}}\par
Code:
\begin{qoptions}
\opt{A}{P $\rightarrow$ 2;Q $\rightarrow$ 3;R $\rightarrow$ 1;S $\rightarrow$ 5}
\opt{B}{P $\rightarrow$ 1;Q $\rightarrow$ 2;R $\rightarrow$ 5;S $\rightarrow$ 1}
\opt{C}{P $\rightarrow$ 2;Q $\rightarrow$ 4;R $\rightarrow$ 3;S $\rightarrow$ 4}
\opt{D}{P $\rightarrow$ 2;Q $\rightarrow$ 2,4;R $\rightarrow$ 4;S $\rightarrow$ 5}
\end{qoptions}
\sect{SECTION-II}
\qn{1} A child has $n$ coins of Rs. 1. Each day he buys exactly one of the three objects either pencil for Rs. 1 or pen for Rs. 2 or book for Rs. 2. If $S_n$ denotes the number of ways of spending $n$ rupees, then the value of $\left(\dfrac{S_{11}-S_{10}}{S_9}+S_7\right)$ is
\qn{2} Let incircle of radius 4 units of a triangle ABC touches the side BC at D. If BD $=6$, DC $=8$ and $\Delta$ be the area of triangle, and $R$ be radius of circumcircle of $\Delta$ABC, then $\sqrt{\Delta-3}+R$ is equal to \underline{\hspace{1.5cm}}.
\qn{3} On a semicircle with diameter AB and centre O point E and C are marked in such a way that OE is perpendicular to AB, and the chord AC intersects the segment OE at a point D which is interior to the semicircle. If $\angle CBA=\dfrac{\pi}{k}$ such that a circle can be inscribed into the quadrilateral OBCD then value of $k$ is equal to
\qn{4} If $z_1,z_2,z_3\in C$ satisfy the system of equations given by $|z_1|=|z_2|=|z_3|=1$, $z_1+z_2+z_3=1$ and $z_1z_2z_3=1$ such that $\mathrm{Im}(z_1)<\mathrm{Im}(z_2)<\mathrm{Im}(z_3)$ then find the value of $\left|z_1+z_2^2+z_3^3\right|$ is \underline{\hspace{1.5cm}}
\qn{5} If $\displaystyle\sum_{n=0}^{\infty}\tan^{-1}\left(\dfrac{\cot^{-1}\left(n^2+3n+3\right)}{1+\cot^{-1}(n+1)\cot^{-1}(n+2)}\right)=\tan^{-1}\left(\dfrac p4\right)$ then find the value of $\cos^{-1}(\cos(p-1))$\par
[Use : $\pi=3.14$]
\qn{6} Given that $f(x)$ is a function defined on R, satisfying $f(1)=1$ and for any $x\in R$, $f(x+5)\ge f(x)+5$ and $f(x+1)\le f(x)+1$. If $g(x)=f(x)+7-x$, then $g(2023)$ is equal to \underline{\hspace{1.5cm}}
\qn{7} Let $f(x)=2x(2-x)$, $0\le x\le2$. If the number of roots and sum of all roots of equation $f(f(f(x)))=\dfrac x2$ are equal to '$a$' and '$b$' respectively, then $(b+a)$ is equal to \underline{\hspace{1.5cm}}
\qn{8} Tangent is drawn at the point $(x_i,y_i)$ on the curve $y=f(x)$, which intersects the x-axis at $(x_{i+1},0)$. Now again a tangent is drawn at $(x_{i+1},y_{i+1})$ on the curve which intersects the x-axis at $(x_{i+2},0)$ and the process is repeated 'n' times. If $x_1,x_2,x_3\ldots.x_n$ form a G.P. with common ratio '$r$' and curve $y=f(x)$ is a hyperbola with eccentricity equal to '$e$' then $r+e^2$ is equal to \underline{\hspace{1.5cm}}
\end{multicols}
\end{document}
